% BEGIN INCLUDED SGA7I_English_current_buildable_source_20260823/components/08_expose_II_complete.tex
%%%%%%%%%%  scan idx 36  |  folio 25  %%%%%%%%%%
\begin{center}
\underline{\SGARefAliasTarget{sga7.sga7i.eii.nroot}报告 II}

\vspace{1em}
\CJKunderdblline{基本群的有限性性质}

\vspace{1em}
\underline{Mich\`ele Raynaud 著}
\end{center}

设 $k$ 是特征为 $p\geq 0$ 的可分闭域，$X$ 是 $k$ 上的连通有限型概形，
$\Pi_1^{(p)}(X)$ 是 $\Pi_1(X)$ 的最大“素于 $p$”商。
本报告将证明：当 $X$ 是曲面时，$\Pi_1^{(p)}(X)$ 在拓扑意义下有限表示；
如果假定奇点解消成立，则对任意 $X$ 也有同样的结论。

对于曲面，J.-P. Murre 的证明方法是归结到一个以 $P^1$ 为底、
具有下述命题中性质 (i)、(ii)、(iii) 的纤维化。

\underline{命题 \SGARefTarget{sga7.sga7i.eii.n2-1}{2.1}。设 $k$ 是代数闭域，
$S$ 是 $k$ 上的射影、不可约、光滑曲面。则存在爆破}
\[
 g:S'\longrightarrow S,
\]
\underline{其中 $S'$ 是正则曲面，且带有纤维化}
\[
 f:S'\longrightarrow P^1,
\]
\underline{使得 $f$ 有一个截面 $\sigma$，并满足下列条件：}

(i) \underline{$f$ 在 $\sigma(P^1)$ 的各点处光滑。}

(ii) \underline{$f$ 的各条纤维都是几何整曲线，且只有通常奇点。}

(iii) \underline{$P^1$ 中存在非空开子集 $U$，使得 $f$ 在 $U$ 各点上的纤维光滑。}

该命题由存在一个容许 Lefschetz 铅笔的射影嵌入
$S\longrightarrow P^r$ 推出（XVII 2.5）。事实上，设
$D=\{H_t\}_{t\in P^1}$ 是这样的超平面铅笔。回顾它的下列性质
（同上，\SGAExternalRef{sga7.sga7i.eii.n2-2}{2.2}{SGA 7 II XVII 2.2}）：
% Source hyperlink erroneously points to the present II 2.2; displayed citation preserved.

%%%%%%%%%%  scan idx 37  |  folio 26  %%%%%%%%%%

(a) $D$ 的轴 $\Delta$ 与 $S$ 横截相交；因此闭子概形 $\Delta\cap S$
是既约的，由 $S$ 的有限个闭点 $x_1,\ldots,x_n$ 组成。

(b) $P^1$ 中存在非空开子集 $U$，使得对每个 $t\in U$，$H_t$ 与 $S$ 横截相交。

(c) 对每个 $t\in P^1-U$，除一个点以外，$H_t$ 与 $S$ 横截相交；
这个例外点是通常二次奇点。

经过 $S$ 中每个不属于 $\Delta$ 的点 $x$，恰有一个超平面 $H_t$。
将 $x$ 对应到 $t\in P^1$ 的映射是有理映射
$a:S\dashrightarrow P^1$，其定义域为 $S-(\Delta\cap S)$。

考虑沿闭子概形 $S\cap\Delta$ 的爆破 $g:S'\longrightarrow S$。
概形 $S'$ 正则，而 $x_i$ 上的纤维 $S'_{x_i}$ 是由 $S$ 在 $x_i$ 处的
切空间定义的射影丛，因而同构于 $P^1$（EGA IV 19.4.2、19.4.4）。
此外，有理映射 $f=ag$ 处处有定义（XVIII 3.1）。最后，每个 $x_i$ 都给出
$f$ 的一个截面：它把 $t\in P^1$ 映到 $H_t\cap S$ 在 $x_i$ 处的切向量。

我们验证条件 (i)、(ii)、(iii)。对每个 $t\in P^1$，曲线 $S'_t$ 与
$S_t=H_t\cap S$ 同构。事实上，态射 $S'_t\longrightarrow S_t$ 除了在
$x_i$ 这些点处可能例外，都是同构，而 $S_t$ 在这些点处正则；此外，
由于 $S'_t$ 是 $S'$ 中的完全交，它是整的，因此同构于 $S_t$。

所断言的纤维性质由 (a)--(c) 推出。最后，$f$ 在 $\sigma$ 的各点处光滑，
因为 $f$ 平坦，且纤维 $S'_t$ 在 $\sigma(t)$ 处几何正则。

%%%%%%%%%%  scan idx 38  |  folio 27  %%%%%%%%%%

\underline{推论 \SGARefTarget{sga7.sga7i.eii.n2-2}{2.2}。沿用
\SGARefLink{sga7.sga7i.eii.n2-1}{2.1} 的记号，设 $Y$ 是 $S$ 中的正规交叉除子，
且 $Y'=Y\times_S S'$。则可选取 $f$、$g$、$U$，使
$Y'\times_{P^1}U$ 为相对于 $U$ 的正规交叉除子。}

已知，如果嵌入 $S\longrightarrow P^r$ 容许 Lefschetz 铅笔，则所有
Lefschetz 铅笔构成 $P^r$ 中直线簇的一个开子集（XVII 3.2.1）。
对于这样的 Lefschetz 铅笔所成的集合，同样如此：它的轴不与 $Y$ 相交，
且 $P^1$ 中存在非空开子集 $U_1$，使得当 $t\in U_1$ 时 $Y\cap H_t$ 光滑。
于是，只需用满足这两个条件的铅笔构造 $f$、$g$，并将 $U$ 换为 $U\cap U_1$。

\SGARefTarget{sga7.sga7i.eii.n2-3-0}{2.3.0}。设 $G$ 是投射有限群。
回顾：若存在整数 $n$，使 $G$ 同构于具有 $n$ 个生成元的自由预群 $F(n)$
的一个商，就称 $G$ 在拓扑意义下有限生成（topologically finitely generated）。
令 $K=\Ker(F(n)\longrightarrow G)$。如果还能在 $K$ 中找到有限个元素
$k_1,\ldots,k_r$，使包含 $k_1,\ldots,k_r$ 的 $K$ 的最小闭正规子群等于 $K$，则称
$G$ 在拓扑意义下有限表示（topologically finitely presented）。
% Review note II-2.3.0: the French original also prints “sous-groupe invariant
% fermé de K”. The ambient normality convention is ambiguous; translation
% preserves the printed wording rather than silently replacing K by F(n).

\underline{定理 \SGARefTarget{sga7.sga7i.eii.n2-3-1}{2.3.1}。设 $k$ 是特征为
$p\geq 0$ 的可分闭域，$X$ 是 $k$ 上的连通有限型概形。假定在 $k$ 的
一个代数闭包上，维数至多为 $\dim X$ 的有限型概形都容许强奇点解消
（SGA 5 I 3.1.5）。由 [1]，当 $\dim X=2$ 时此条件成立；由 [2]，当
$\operatorname{char}k=0$ 时也成立。则基本群 $\Pi_1^{(p)}(X)$ 在拓扑意义下有限表示。}

由 SGA 1 IX 4.10，可以假定 $k$ 代数闭。

1) \underline{归结到 $X$ 是正则曲面的情形。}

若 $f:X'\longrightarrow X$ 是有限型态射，令 $X''=X'\times_X X'$，
并分别用 $(X'_a)_{a\in A}$ 与 $(X''_b)_{b\in B}$ 表示 $X'$、$X''$ 的连通分支族。
当 $f$ 是有限平展覆盖范畴的有效下降态射时，基本群 $\Pi_1(X)$
可以用基本群 $\Pi_1(X'_a)$ 和 $\Pi_1(X''_b)$ 显式地算出。

%%%%%%%%%%  scan idx 39  |  folio 28  %%%%%%%%%%

由此可知（SGA 1 IX 5.2、\SGAExternalRef{sga7.sga7i.eix.n5-3}{5.3}{SGA 1 IX 5.3}），
若各 $\Pi_1(X'_a)$ 在拓扑意义下有限生成，则 $\Pi_1(X)$ 也如此；若各
$\Pi_1(X'_a)$ 在拓扑意义下有限表示，而各 $\Pi_1(X''_b)$ 在拓扑意义下
有限生成，则 $\Pi_1(X)$ 在拓扑意义下有限表示。

由于 $X$ 容许奇点解消，存在正则概形 $X'$ 及紧合双有理态射
$f:X'\longrightarrow X$。因 $f$ 是有限平展覆盖的有效下降态射，只需在
附加假设 $X$ 正则的条件下证明定理。首先推出 $\Pi_1(X)$ 在拓扑意义下
有限生成；再利用 $\Pi_1(X')$ 在拓扑意义下有限表示、各 $\Pi_1(X''_b)$
在拓扑意义下有限生成，便知 $\Pi_1(X)$ 在拓扑意义下有限表示。

现取 $X$ 的有限仿射开覆盖 $(U_i)_{i\in I}$，令
$f:\bigsqcup U_i\longrightarrow X$ 为典范态射。由于 $f$ 是有限平展覆盖的
有效下降态射，只需对各 $U_i$ 证明定理。因此归结到 $X$ 仿射且光滑的情形。

此后假定 $X$ 在 $k$ 上仿射且光滑，且 $\dim X>2$。
可应用拟射影概形的 Lefschetz 定理（[3] V 7.4）：若 $Y$ 是 $X$ 的足够一般的
超平面截面，则
\[
 \Pi_1(Y)\simeq\Pi_1(X).
\]
因而，只需对 $Y$ 证明定理，这就归结到 $\dim X=2$ 的情形。

2) \underline{$X$ 是 $k$ 上光滑曲面的情形。}

由曲面的强奇点解消定理，存在光滑、整、射影曲面 $S$ 及开浸入
$i:X\longrightarrow S$，使得 $Y=S-X$ 是正规交叉除子。
应用推论 \SGARefLink{sga7.sga7i.eii.n2-2}{2.2}，并沿用其中的记号。

%%%%%%%%%%  scan idx 40  |  folio 29  %%%%%%%%%%

我们证明，只需证明 $X'=X\times_S S'$ 的基本群在拓扑意义下有限表示。
由于 $X$、$X'$ 正则，且
\[
 \operatorname{codim}\!\left(\bigcup_{1\leq i\leq n}\{x_i\},X\right)
 \geq 2,
\]
$X'$ 的每个平展覆盖都在 $X-\bigcup_{1\leq i\leq n}\{x_i\}$ 上诱导一个
平展覆盖，而后者唯一延拓为 $X$ 的平展覆盖（SGA 2 1.11）。这给出同构
\[
 \Pi_1(X')\simeq\Pi_1(X).
\]

考虑笛卡儿图
\[
\begin{tikzcd}[column sep=large,row sep=large]
X' \arrow[d,"h"'] & X'_U \arrow[l] \arrow[d,"h_U"] \\
P^1 & U \arrow[l]
\end{tikzcd}
\]
其中 $h=f|_{X'}$。设 $L$ 为异于 $p$ 的素数所成的集合，$\bar\eta$ 为
$P^1$ 的一般点 $\eta$ 上的几何点。我们验证 SGA 1 XIII 4.8 的假设。
态射 $h$ 有截面 $\sigma$，因为 $f$ 有这样的截面。为证明 SGA 1 XIII 4.7 的
条件 (a)，先注意到，由于 $h$ 的纤维几何整，$h$ 是 $0$-无环且局部
$0$-无环的（SGA 4 XV 4.1、1.16）。还须说明：对 $X'$ 上的任意局部常值
集合层 $F$，$(F,h)$ 在维数 $\leq 0$ 时上同调紧合。若 $\xi$ 为 $P^1$ 的
闭点，并用 $Z$ 表示 $\Spec\mathcal O_{P^1,\xi}$ 在 $\bar\eta$ 中的整闭包，
则 $X'_Z=X'\times_{P^1}Z$ 在其闭点处满足条件 $(S_2)$，在其余各点处正则，
从而正规。因此，$X'_Z$ 的一个平展覆盖，若其在 $X'_{\bar\eta}$ 上的限制
连通，本身就必然连通。这证明了所断言的结论，因而也证明了条件 (a)。

由于 $X'_U$ 是 $S'_U$ 中一个相对于 $U$ 的正规交叉除子的补集，$h_U$
对 $L$ 局部 $1$-非球面（SGA 4 XV 2.1）；且对 $X'_U$ 上的每个有限常值
$L$-群层 $F$，$(F,h_U)$ 在维数 $\leq 1$ 时上同调紧合
（SGA 1 XIII 2.4），这给出条件 (b)。由于 $P^1$ 正规，且
$T=P^1-U$ 的余维为 $1$，有

%%%%%%%%%%  scan idx 41  |  folio 30  %%%%%%%%%%

\[
 \operatorname{prof\,\acute et}_T(P^1)\geq 2,
\]
这给出条件 (c)。由纯性定理，条件
\[
 \operatorname{prof\,hop}_x(X')\geq 3
\]
在 $X'$ 的每个闭点处成立，这给出条件 (e)。最后，对每个 $t\in P^1$，
$h$ 在 $\sigma(t)$ 处光滑，从而在 $\sigma(t)$ 处对 $L$ 局部 $1$-非球面，
这给出条件 (d)。

由 SGA 1 XIII 4.7，若 $a$ 为 $X'_{\bar\eta}$ 的几何点，则有正合列
\[
 1\longrightarrow\Pi_1^{(p)}(X'_{\bar\eta},a)
 \longrightarrow\Pi'_1(X'_U,a)
 \longrightarrow\Pi_1(U,a)\longrightarrow 1,
\]
而选定截面 $\sigma$ 就定义了此正合列的一个分裂。
设 $K$ 是 $\Pi_1(U,a)$ 在 $\Aut(\Pi_1^{(p)}(X'_{\bar\eta},a))$ 中的像。
在 $K$ 作用下的余不变量群
\[
 \Pi_1^{(p)}(X'_{\bar\eta},a)_K,
\]
是 $\Pi_1^{(p)}(X'_{\bar\eta},a)$ 除以由形如 $x^{-1}q(x)$ 的元素生成的
闭正规子群所得的商，其中 $x\in\Pi_1^{(p)}(X'_{\bar\eta},a)$，$q\in K$。
由同上 (4.7.4)，有同构
\[
 \Pi_1^{(p)}(X',a)\simeq
 \Pi_1^{(p)}(X'_{\bar\eta},a)_K.
\]

写成 $P^1-U=\coprod_{1\leq i\leq n}\{t_i\}$。对每个 $i$，令
$\bar T_i$ 为 $P^1$ 在 $t_i$ 上的几何点 $\bar t_i$ 处的严格局部化，
令 $s_i$ 为 $\bar T_i$ 的一般点。$U$ 的一个平展覆盖，若在每个 $s_i$ 上的
限制平凡，就延拓为 $P^1$ 的平展覆盖，从而自身平凡。对每个 $i$，令
$\overline{k(s_i)}$ 为 $k(s_i)$ 的代数闭包。前面的断言说明，$\Pi_1(U,a)$
是包含各 Galois 群 $\Gal(\overline{k(s_i)}/k(s_i))$ 在 $\Pi_1(U,a)$ 中的像的
最小闭正规子群。由 I \SGARefLink{sga7.sga7i.ei.n5-3}{5.3}，
$\Gal(\overline{k(s_i)}/k(s_i))$ 在 $\Aut(\Pi_1^{(p)}(X'_{\bar\eta},a))$ 中的像
$K_i$ 有限。由于 $K$ 是包含所有 $K_i$ 的最小闭正规子群，且
$\Pi_1^{(p)}(X'_{\bar\eta},a)$ 在拓扑意义下有限表示（SGA 1 XIII 3.2），
这证明 $\Pi_1^{(p)}(X'_{\bar\eta},a)_K$，从而与其同构的群
$\Pi_1^{(p)}(X)$，在拓扑意义下有限表示。

%%%%%%%%%%  scan idx 42  |  folio 31  %%%%%%%%%%

\underline{注 \SGARefTarget{sga7.sga7i.eii.n2-3-2}{2.3.2}。}若 $X$ 是一个
射影且光滑的 $k$-概形中某个正规交叉除子的开补集，则证明定理
\SGARefLink{sga7.sga7i.eii.n2-3-1}{2.3.1} 无须使用奇点解消。

\begin{center}\rule{1.5cm}{0.4pt}\end{center}
\begin{center}\underline{参考文献}\end{center}

\begin{enumerate}[label={[\arabic*]},leftmargin=3em]
\item S. Abhyankar, ``Local uniformization on algebraic surfaces over
ground fields of characteristic $p\neq0$,'' \emph{Ann. of Math.},
第 63 卷，第 3 期，1956 年，第 491--526 页。

\item H. Hironaka, ``Resolution of singularities of an algebraic variety
over a field of characteristic zero,'' \emph{Ann. of Math.}, 第 79 卷，
1964 年。

\item M. Raynaud, \emph{Th\'eor\`eme de Lefschetz en cohomologie des
faisceaux coh\'erents et en cohomologie \'etale}，学位论文，即将出版。
\end{enumerate}
% END INCLUDED SGA7I_English_current_buildable_source_20260823/components/08_expose_II_complete.tex
